\section{Ključni izvori rada}

Kao dva ključna izvora za ovaj seminarski rad izdvajaju se rad Pabla Gonzalesa de la Torea (\textit{Pablo González de la Torre}), Marte Peres-Verdugo (\textit{Marta Pérez-Verdugo}) i Šabijera E. Barandijarana (\textit{Xabier E. Barandiaran}) \textit{Attention is all they need: Cognitive science and the (techno)political economy of attention in humans and machines}, kao i rad Atuse Kasirzade (\textit{Atoosa Kasirzadeh}) i Čarlsa Evansa (\textit{Charles Evans}) \textit{User Tampering in Reinforcement Learning Recommender Systems}.

Prvi rad je ključan za širi teorijski okvir seminarskog rada. U njemu se pažnja ne posmatra samo kao individualna psihološka sposobnost, već i kao resurs koji savremeni digitalni sistemi mere, usmeravaju i ekonomski koriste. Posebno je važan zato što povezuje kognitivno razumevanje pažnje sa političkom ekonomijom digitalnih platformi i ukazuje na to da se pažnja u digitalnom okruženju oblikuje i kroz šire društvene i tehnološke procese, a ne samo kroz individualne navike korisnika. Na osnovu ovog izvora se u daljem tekstu razvija okvir ekonomije pažnje, odnosno ideja da se pažnja korisnika u digitalnom okruženju pretvara u merljivu i poslovno upotrebljivu vrednost~\cite{gonzalez2024attention}.

Drugi rad je ključan za tehničko-etički deo seminarskog rada. Kasirzade i Evans analiziraju problem manipulacije korisnikom (engl. \textit{user tampering}) u sistemima preporuke zasnovanim na učenju potkrepljivanjem. Njihov rad je značajan zato što pokazuje da sistem preporuke ne mora samo da prati postojeće preferencije korisnika, već može naučiti da ih oblikuje ukoliko mu to pomaže da dugoročno poveća korisničko angažovanje. Posebno je zanimljiv njihov simulacioni primer političkog sistema preporuke, u kome se pokazuje kako rane preporuke mogu voditi ka polarizaciji korisnika ako takva promena kasnije olakšava optimizaciju cilja sistema~\cite{kasirzadeh2021user}. Time se jasno otvara pitanje granice između nenametljivog predlaganja sadržaja i ciljane manipulacije korisnikom, naročito onda kada su ciljne promenljive sistema preporuke definisane kroz dugoročno povećanje angažovanja.

Zajedno, ova dva rada povezuju teorijski i tehnički deo seminarskog rada. Prvi objašnjava zašto je pažnja postala centralni resurs digitalnog kapitalizma, dok drugi pokazuje kako konkretni algoritamski sistemi mogu učestvovati u oblikovanju ponašanja korisnika. Zbog toga oni predstavljaju osnovu za glavnu tezu rada: da problem korisničkog angažovanja nije samo pitanje merenja popularnosti sadržaja, već i pitanje načina na koji se pažnja predviđa, optimizuje i potencijalno manipuliše.
